\section{Implementation details}
\begin{multicols}{2}
\todo[inline, color=green!40]{描述training, testing data的選取}

The training data of each session are the first 5000 time bins, and the rest is the testing data.
Figure \ref{fig:full_testing_data} visualized the complete testing data from a few sessions,
whereas some sessions, had their testing data trimmed by \cite{makin2018} due to the abnormal recording periods.
These sessions contained intervals when monkeys fingertip stayed still for a considerable long while.
As marked in Figure \ref{fig:trimmed_testing_data}, shaded red parts indicated the discarded intervals.
\end{multicols}

\begin{figure}[H]
    \foreach \i in {indy_20160622_01,indy_20160624_03,indy_20160627_01}
        {
            \begin{subfigure}{1\textwidth}
                \includegraphics[height=0.09\textheight]{../Figures_2/\i_x_y_testing_data_trajectory.png}
            \end{subfigure}
        }
        \captionsetup{width=0.5\textwidth}
        \caption{Testing Data with Full Length}
        \label{fig:full_testing_data}
\end{figure}

\newpage
\todo[author=Wu, color=blue!40]{10個session剛好一整頁}
\begin{figure}[H]
    \foreach \i in {loco_20170210_03,loco_20170213_02,loco_20170214_02,loco_20170215_02,loco_20170216_02,loco_20170217_02,loco_20170227_04,loco_20170228_02,loco_20170301_05,loco_20170302_02}
        {
            \begin{subfigure}{1\textwidth}
                \includegraphics[height=0.09\textheight]{../Figures_2/\i_x_y_testing_data_trajectory.png}
            \end{subfigure}
        }
        \captionsetup{width=0.5\textwidth}
        \caption{Trimmed Testing Data}
        \label{fig:trimmed_testing_data}
\end{figure}


\section{Qualitative analysis}
\label{sec:qualitative_analysis}
\begin{multicols}{2}
\todo[inline, color=green!40]{放入 Figure~\ref{fig:position_prediction} ~ Figure~\ref{fig:acceleration_prediction}，將縱軸統一}
\todo[inline, color=green!40]{ bidirectional GRU + L.N. + attention  結果，分別畫出位置、速度、加速度對時間的預測圖。Position 的解碼結果可以畫出 2D 平面軌跡圖。 }
This section we show the result of our predicted variables. 
Each figure represents a session, and here we selected 5 sessions for one type of variable.
Since our models were designed to decode both x and y axis of a variable at the same time, 
it is convenient to plot the predictions of the two axes in the same figure, 
also easiler for us to compare their results.

\end{multicols}

\newpage
\todo[inline, author=Wu, color=blue!40]{ 位置、速度、加速度 各一頁 }
\begin{figure}[H]
    \foreach \i in {1,2,3,4,5}
    {
        \begin{subfigure}{1\textwidth}
            \includegraphics[height=0.18\textheight]{../KM_reconstruction_plots/\i.png}
        \end{subfigure}
    }
    \caption{Position Predictions}
    \label{fig:position_prediction}
\end{figure}

\newpage
\todo[inline, author=Wu, color=blue!40]{ 位置、速度、加速度 各一頁 }
\begin{figure}[H]
    \foreach \i in {6,7,8,9,10}
    {
        \begin{subfigure}{1\textwidth}
            \includegraphics[height=0.18\textheight]{../KM_reconstruction_plots/\i.png}
        \end{subfigure}
    }
    \caption{Velocity Predictions}
    \label{fig:velocity_prediction}
\end{figure}

\newpage
\todo[inline, author=Wu, color=blue!40]{ 位置、速度、加速度 各一頁 }
\begin{figure}[H]
    \foreach \i in {11,12,13,14,15}
    {
        \begin{subfigure}{1\textwidth}
            \includegraphics[height=0.18\textheight]{../KM_reconstruction_plots/\i.png}
        \end{subfigure}
    }
    \caption{Acceleration Predictions}
    \label{fig:acceleration_prediction}
\end{figure}


\begin{figure}[H]
    \begin{center}
        \includegraphics[width=1\textwidth]{Firing_rate_heatmap_and_all_kinematic_variables_pred_and_actual_no_color_bar.png}
        \captionsetup{width=0.5\textwidth}
        \caption{Firing rate and decoding result}
        \label{fig:Firing_rate_heatmap_and_all_kinematic_variables_pred_and_actual_no_color_bar}
    \end{center}
\end{figure}



\section{Comparison of state-of-the-art}
\begin{multicols}{2}
\todo[inline, color=green!40]{包含Wiener filter、KF、簡單的neural network、Makin的方法，每個方法引用一個經典的論文，放入以下論文的Figure 4
Superior arm-movement decoding from cortex with a new, unsupervised-learning algorithm}
\end{multicols}


\section{Visualization of attention process}
\begin{multicols}{2}
\todo[inline, color=green!40]{放入attention map}
We have implemented the attention mechanism to both one-way and bidirectional GRUs.
To maximize the information we can get form the attention map visualization,
we put the predicted trajectories, attention map and the input neural firing rate torgether as showed in Figure~\ref{fig:decoding_trajectories_with_attention}.
\end{multicols}

\begin{figure}[H]
    \begin{center}
        \begin{subfigure}{0.6\textwidth}
            \includegraphics[width=1\textwidth]{bidir_pos_atten.png}
            \caption{ Position }
        \end{subfigure}
        \begin{subfigure}{0.6\textwidth}
            \includegraphics[width=1\textwidth]{bidir_vel_atten.png}
            \caption{ Velocity }
        \end{subfigure}
        \begin{subfigure}{0.6\textwidth}
            \includegraphics[width=1\textwidth]{bidir_acc_atten.png}
            \caption{ Acceleration }
        \end{subfigure}
        \caption{ Trajectories with attention map }
        \label{fig:decoding_trajectories_with_attention}
    \end{center}
\end{figure}


% \begin{landscape}
    \section{Ablation study}
    \subsection{Comparison table}
    \label{sec:ablation_study}
    \begin{table}[H]
        \centering
        \caption{Weighted average $R^2$ and standard deviation}
        \begin{tabular}{ccccccccc}
            \hline
            \multicolumn{1}{c}{ \multirow{2}{*}{\textbf{Atten.}}}  & \multicolumn{1}{c}{\multirow{2}{*}{\textbf{L.N.}}}    & \multicolumn{1}{c|}{\multirow{1}{*}{\textbf{GRU}}}          & \multicolumn{2}{c}{\multirow{1}{*}{\textbf{Position}}} & \multicolumn{2}{c}{\multirow{1}{*}{\textbf{Velocity}}} & \multicolumn{2}{c}{\multirow{1}{*}{\textbf{Acceleration}}} \\
            \multicolumn{1}{c}{}                                   & \multicolumn{1}{c}{}            & \multicolumn{1}{c|}{\multirow{1}{*}{\textbf{Dir}}} & \multicolumn{1}{c}{\multirow{1}{*}{\textbf{x}}} & \multicolumn{1}{c}{\multirow{1}{*}{\textbf{y}}} & \multicolumn{1}{c}{\multirow{1}{*}{\textbf{x}}} & \multicolumn{1}{c}{\multirow{1}{*}{\textbf{y}}} & \multicolumn{1}{c}{\multirow{1}{*}{\textbf{x}}} & \multicolumn{1}{c}{\multirow{1}{*}{\textbf{y}}} \\
            \hline
            \textcolor{red}{\xmark} & \textcolor{red}{\xmark} & $\rightarrow$                   & 0.49$\pm$0.21         & 0.68$\pm$0.12         & 0.65$\pm$0.13         & 0.71$\pm$0.11         & 0.31$\pm$0.16         & 0.32$\pm$0.23 \\ 
            \textcolor{red}{\xmark} & \textcolor{red}{\xmark} & $\rightleftarrows$              & 0.49$\pm$0.21         & 0.69$\pm$0.13         & 0.65$\pm$0.12         & 0.74$\pm$0.07         & 0.24$\pm$0.13         & 0.25$\pm$0.19 \\ 
            \textcolor{red}{\xmark} & \textcolor{blue}{\checkmark} & $\rightarrow$              & 0.53$\pm$0.19         & 0.69$\pm$0.09         & 0.64$\pm$0.11         & 0.75$\pm$0.05         & 0.40$\pm$0.13         & 0.47$\pm$0.10 \\ 
            \textcolor{red}{\xmark} & \textcolor{blue}{\checkmark} & $\rightleftarrows$         &\textbf{0.55$\pm$0.17} & 0.70$\pm$0.08         & 0.53$\pm$0.19         & 0.71$\pm$0.07         &\textbf{0.41$\pm$0.13} & 0.50$\pm$0.09 \\ 
            \textcolor{blue}{\checkmark} & \textcolor{red}{\xmark} & $\rightarrow$              & 0.46$\pm$0.22         & 0.67$\pm$0.11         & 0.62$\pm$0.12         & 0.73$\pm$0.06         & 0.12$\pm$0.11         & 0.14$\pm$0.15 \\
            \textcolor{blue}{\checkmark} & \textcolor{red}{\xmark} & $\rightleftarrows$         & 0.49$\pm$0.21         & 0.69$\pm$0.09         & 0.65$\pm$0.11         & 0.74$\pm$0.05         & 0.25$\pm$0.19         & 0.29$\pm$0.22 \\
            \textcolor{blue}{\checkmark}  & \textcolor{blue}{\checkmark} & $\rightarrow$        & 0.51$\pm$0.20         & 0.69$\pm$0.09         & 0.62$\pm$0.12         & 0.74$\pm$0.05         & 0.35$\pm$0.13         & 0.43$\pm$0.11 \\ 
            \textcolor{blue}{\checkmark} & \textcolor{blue}{\checkmark} & $\rightleftarrows$    & 0.53$\pm$0.19         &\textbf{0.71$\pm$0.08} &\textbf{0.66$\pm$0.10} &\textbf{0.76$\pm$0.05} & 0.40$\pm$0.13         &\textbf{0.48$\pm$0.09} \\ 
            \hline
            \multicolumn{3}{c}{ \cite{makin2018} } & 0.46$\pm$0.28 & 0.70$\pm$0.09  & 0.58$\pm$0.15 & 0.74$\pm$0.06 & 0.29$\pm$0.15 & 0.40$\pm$0.09 \\
            \hline
        \end{tabular}
        \label{tab:ablation_study}
    \end{table}
    % \end{landscape}
    
    
    \subsection{With or without attention in bidirectional GRU}
    
    \begin{multicols}{2}
    There are some sessions does not benefit from our attention mechanism.
    Their prediction $R^2$ decreased when comparing to the result from the decoders without implementing attention.
    Take session \textbf{indy\_20161027\_03} for example,
    the $R^2$ in position decoding from the bidirectional LN-GRU is $0.64$, $0.66$ in x-axis and y-axis.
    But the result from the same decoder with attention, the $R^2$ becomes $0.44$ and $0.58$, 
    with a significant drop in x-axis position prediction.
    Figure~\ref{fig:indy_20161027_03_position_x_fauilure_case_comparison} compares the position decoding results along x-axis.
    Both decoders tend to fail when the movement of the fingertip stopped.
    However in this session, decoder with attention gives more deviation to the predicted trajectory than the decoder without attention.
    \end{multicols}
    
    \begin{figure}[H]
        \begin{center}
            \begin{subfigure}{1\textwidth}
                \includegraphics[height=0.15\textheight]{../Failure_case_figures/indy_20161027_03/x_pos_vs_time_from_1368_to_1381.png}
                \caption{}
            \end{subfigure}
            \begin{subfigure}{1\textwidth}
                \includegraphics[height=0.15\textheight]{../Failure_case_figures/indy_20161027_03/x_pos_vs_time_from_1509_to_1522.png}
                \caption{}
            \end{subfigure}
            \caption{x position prediction in session indy\_20161027\_03}
            \label{fig:indy_20161027_03_position_x_fauilure_case_comparison}
        \end{center}
    \end{figure}
    
    Figure~\ref{fig:indy_20161027_03_position_y_fauilure_case_comparison} shows that the predicted trajectory in y-axis is smoother than the one without attention.
    But in the intervals with the fingertip stayed stationary, as from $1513s$ to $1516s$,
    the deviation is larger the the trajectory predicted without attention.
    
    
    \begin{figure}[H]
        \begin{center}
            \begin{subfigure}{1\textwidth}
                \includegraphics[height=0.15\textheight]{../Failure_case_figures/indy_20161027_03/y_pos_vs_time_from_1368_to_1381.png}
                \caption{}
            \end{subfigure}
            \begin{subfigure}{1\textwidth}
                \includegraphics[height=0.15\textheight]{../Failure_case_figures/indy_20161027_03/y_pos_vs_time_from_1509_to_1522.png}
                \caption{}
            \end{subfigure}
            \caption{y position prediction in session indy\_20161027\_03}
            \label{fig:indy_20161027_03_position_y_fauilure_case_comparison}
        \end{center}
    \end{figure}
    
    \begin{multicols}{2}
    The corresponding attention map is showed in Figure~\ref{fig:indy_20161027_03_corresponding_attention_map},
    and we can tell that our decoder with attention failed when it has lose the focus a the time bins.
    In other words, the attention scores are distributed rather evenly among all the past time bins,
    whereas in the samples of good prediction accuracy, 
    the attention scores are more concentrated on some specific past time bins.
    our attention should able to distinguish if the fingertip of the subjects is moving or stationary.
    This not an easy task since neurons still fire when the subjects hold their movement.
    
    \end{multicols}
    
    \begin{figure}[H]
        \begin{center}
            \begin{subfigure}{0.49\textwidth}
                \includegraphics[width=1\textwidth]{../Failure_case_figures/indy_20161027_03/attention_map_600_to_800.png}
                \caption{}
            \end{subfigure}
            \begin{subfigure}{0.49\textwidth}
                \includegraphics[width=1\textwidth]{../Failure_case_figures/indy_20161027_03/attention_map_2800_to_3000.png}
                \caption{}
            \end{subfigure}
            \caption{Corresponding attention map in session indy\_20161027\_03}
            \label{fig:indy_20161027_03_corresponding_attention_map}
        \end{center}
    \end{figure}
    
    \subsection{One-way or bidirectional GRU}
    \begin{multicols}{2}
    
    \end{multicols}
    
    \begin{figure}[H]
        \begin{center}
            \begin{subfigure}{0.49\textwidth}
                \includegraphics[width=1\textwidth]{oneway_pos_atten.png}
                \caption{one-way GRU}
            \end{subfigure}
            \begin{subfigure}{0.49\textwidth}
                \includegraphics[width=1\textwidth]{bidir_pos_atten.png}
                \caption{bidirectional GRU}
            \end{subfigure}
            \caption{Position decoder with attention}
            \label{fig:position_decoder_with_attention}
        \end{center}
    \end{figure}
    
    \begin{figure}[H]
        \begin{center}
            \begin{subfigure}{0.49\textwidth}
                \includegraphics[width=1\textwidth]{oneway_vel_atten.png}
                \caption{one-way GRU}
            \end{subfigure}
            \begin{subfigure}{0.49\textwidth}
                \includegraphics[width=1\textwidth]{bidir_vel_atten.png}
                \caption{bidirectional GRU}
            \end{subfigure}
            \caption{Velocity decoder with attention}
            \label{fig:velocity_decoder_with_attention}
        \end{center}
    \end{figure}
    
    \begin{figure}[H]
        \begin{center}
            \begin{subfigure}{0.49\textwidth}
                \includegraphics[width=1\textwidth]{oneway_acc_atten.png}
                \caption{one-way GRU}
            \end{subfigure}
            \begin{subfigure}{0.49\textwidth}
                \includegraphics[width=1\textwidth]{bidir_acc_atten.png}
                \caption{bidirectional GRU}
            \end{subfigure}
            \caption{Acceleration decoder with attention}
            \label{fig:acceleration_decoder_with_attention}
        \end{center}
    \end{figure}

\section{Internal dynamics of network}
\label{sec:internal_dynamics_of_network}
\begin{multicols}{2}
\todo[inline, color=green!40]{參考以下論文的Figure 7 (ClosedLoopLSTM2019)Decoding Movements from Cortical Ensemble Activity Using a Long Short-Term Memory Recurrent Network}
\end{multicols}


\begin{figure}[H]
    \begin{center}
        \includegraphics[width=0.7\textwidth]{internal_dynamic_pos.png}
        \captionsetup{width=0.5\textwidth}
        \caption{ GRU internal dynamics in position decoder }
        \subcaption*{session indy\_20160411\_02}
        \label{fig:internal_dynamic_pos}
    \end{center}
\end{figure}

\begin{figure}[H]
    \begin{center}
        \includegraphics[width=0.7\textwidth]{internal_dynamic_vel.png}
        \captionsetup{width=0.5\textwidth}
        \caption{ GRU internal dynamics in velocity decoder }
        \subcaption*{session indy\_20160407\_02}
        \label{fig:internal_dynamic_vel}
    \end{center}
\end{figure}

\begin{figure}[H]
    \begin{center}
        \includegraphics[width=0.7\textwidth]{internal_dynamic_acc.png}
        \captionsetup{width=0.5\textwidth}
        \caption{ GRU internal dynamics in acceleration decoder }
        \subcaption*{session indy\_20160411\_01}
        \label{fig:internal_dynamic_vel}
    \end{center}
\end{figure}

\section{Failure case study}
\todo[inline, color=green!40]{放入失敗的例子}

\section{Decodeing results from dataset Chewie}
\begin{multicols}{2}
\todo[inline, color=green!40]{ 猴子C的Decoding 結果單獨一個子章節顯示，並不列入與 Makin 的比較 }
\end{multicols}


% \section{y-axis outperforms x-axis}
% \begin{multicols}{2}
% \todo[inline]{用 Mutual Information 解釋 y 方向解碼比 x 方向 解碼還要好的現象}
% Where we used the tool from \cite{ince2017statistical} named  Gaussian-copula mutual information (GCMI) to analysis the reason why in three kinds of kinematic variables,
% the prediction of y-axis always exceed x-axis.
% The module of this tool is available on \url{https://github.com/robince/gcmi}.
% \end{multicols}

